%%%%%%%%%%%%%%%%
%%% gen 9 %%%
%%%%%%%%%%%%%%%%

\Chapter{9}

\Verse{1}
{
	\hP{וַיְבָרֶךְ אֱלֹהִים אֶת נֹחַ וְאֶת בָּנָיו וַיֹּאמֶר לָהֶם פְּרוּ וּרְבוּ וּמִלְאוּ אֶת הָאָרֶץ׃}
}
{
	\eP{
		9 And God blessed Noah and his sons and said to them, “Be\\
		fruitful and multiply and fill the earth.\\
	}
}
{
}

\Verse{2}
{
	\hP{וּמוֹרַאֲכֶם וְחִתְּכֶם יִהְיֶה עַל כּל חַיַּת הָאָרֶץ וְעַל כּל עוֹף הַשָּׁמָיִם בְּכֹל אֲשֶׁר תִּרְמֹשׂ הָאֲדָמָה וּבְכל דְּגֵי הַיָּם בְּיֶדְכֶם נִתָּנוּ׃}
}
{
	\eP{
		And fear of you and dread of you will be on every living\\
		thing of the earth and on every bird of the skies, in\\
		every one that will creep on the earth and in all the fish\\
		of the sea. They are given into your hand.\\
	}
}
{
}

\Verse{3}
{
	\hP{כּל רֶמֶשׂ אֲשֶׁר הוּא חַי לָכֶם יִהְיֶה לְאכְלָה כְּיֶרֶק עֵשֶׂב נָתַתִּי לָכֶם אֶת כֹּל׃}
}
{
	\eP{
		Every creeping every animal that is alive will be yours\\
		for food; I’ve given every one to you like a plant of\\
		vegetation,\\
	}
}
{
%	\footnote{
%		Every creeping animal. Elsewhere this term refers only to certain
%		animals that crawl on the ground, but here it is understood to refer to
%		all animals: “I’ve given one to you like a plant”—meaning that humans
%		are permitted to eat animals in addition to plants. Do humans eat
%		animals before this? Noah makes burnt offerings (lt) after the flood,
%		but that is a type of sacrifice that is not eaten. Abel brings a
%		different kind of offering (minh) to God, which is a type of sacrifice
%		that is understood later in the Torah to be eaten by humans. But, prior
%		to Genesis 9, all statements from the deity to humans concerning eating
%		refer only to plants (1:29; 2:16; 3:18). It is therefore commonly
%		understood that humans are not permitted to eat animals until after the
%		flood. Some say that God now permits it as a concession to human
%		appetites. I see no grounds in the text for this view. A more likely
%		explanation for the change in divine requirements is that humans have
%		saved all animal species by taking them in the ark. All animals now owe
%		their lives to humankind, and this may be why humans are now permitted
%		to take their lives for food.
%	}
}

\Verse{4}
{
	\hP{אַךְ בָּשָׂר בְּנַפְשׁוֹ דָמוֹ לֹא תֹאכֵלוּ׃}
}
{
	\eP{except you shall not eat flesh in its life, its blood,}
}
{
}

\Verse{5}
{
	\hP{וְאַךְ אֶת דִּמְכֶם לְנַפְשֹׁתֵיכֶם אֶדְרֹשׁ מִיַּד כּל חַיָּה אֶדְרְשֶׁנּוּ וּמִיַּד הָאָדָם מִיַּד אִישׁ אָחִיו אֶדְרֹשׁ אֶת נֶפֶשׁ הָאָדָם׃}
}
{
	\eP{
		and except I shall inquire for your blood, for your lives.\\
		I shall inquire for it from the hand of every animal and\\
		from the hand of a human. I shall inquire for a human’s\\
		life from the hand of each man for his brother.\\
	}
}
{
%	\footnote{
%		I shall inquire for your blood, for your lives. This expression means
%		that it is forbidden to take a human life unlawfully. It does not
%		necessarily mean that one may never take a human life; for example, it
%		may not forbid taking a life in self-defense. Rather, God will
%		inquire—that is, require an accounting—and hold the guilty responsible,
%		whether human or animal.
%	}
}

\Verse{6}
{
	\hP{שֹׁפֵךְ דַּם הָאָדָם בָּאָדָם דָּמוֹ יִשָּׁפֵךְ כִּי בְּצֶלֶם אֱלֹהִים עָשָׂה אֶת הָאָדָם׃}
}
{
	\eP{
		One who sheds a human’s blood: by a human his blood will\\
		be shed, because He made the human in the image of God.\\
	}
}
{
}

\Verse{7}
{
	\hP{וְאַתֶּם פְּרוּ וּרְבוּ שִׁרְצוּ בָאָרֶץ וּרְבוּ בָהּ׃}
}
{
	\eP{
		And you, be fruitful and multiply, swarm in the earth and\\
		multiply in it.”\\
	}
}
{
}

\Verse{8}
{
	\hP{וַיֹּאמֶר אֱלֹהִים אֶל נֹחַ וְאֶל בָּנָיו אִתּוֹ לֵאמֹר׃}
}
{
	\eP{And God said to Noah and to his sons with him, saying,}
}
{
}

\Verse{9}
{
	\hP{וַאֲנִי הִנְנִי מֵקִים אֶת בְּרִיתִי אִתְּכֶם וְאֶת זַרְעֲכֶם אַחֲרֵיכֶם׃}
}
{
	\eP{
		“And I: here, I’m establishing my covenant with you and\\
		with your seed after you\\
	}
}
{
%	\footnote{
%		covenant. The flood event concludes with a contract between the deity
%		and humans: a covenant. The Noahic covenant is the first of four major
%		covenants that provide the structure in which nearly all of the Bible is
%		framed. (The exceptions, most notably the book of Job, are relatively
%		few.) The Noahic covenant promises the security of the cosmos. The
%		Abrahamic covenant (Genesis 15 and 17) promises the land to Abraham’s
%		descendants and makes YHWH their God. The Israelite (or Sinai, or
%		Mosaic) covenant (Exodus 20; 34; Deuteronomy 5; 7:12–15) promises
%		well-being in the promised land. And the Davidic covenant (2 Samuel 7;
%		Psalms 89; 132) promises kingship over Jerusalem and Judah to David and
%		his descendants.
%	}
}

\Verse{10}
{
	\hP{וְאֵת כּל נֶפֶשׁ הַחַיָּה אֲשֶׁר אִתְּכֶם בָּעוֹף בַּבְּהֵמָה וּבְכל חַיַּת הָאָרֶץ אִתְּכֶם מִכֹּל יֹצְאֵי הַתֵּבָה לְכֹל חַיַּת הָאָרֶץ׃}
}
{
	\eP{
		and with every living being that is with you, of the\\
		birds, of the domestic animals, and of all the wild\\
		animals of the earth with you, from all those coming out\\
		of the ark to every living thing of the earth.\\
	}
}
{
}

\Verse{11}
{
	\hP{וַהֲקִמֹתִי אֶת בְּרִיתִי אִתְּכֶם וְלֹא יִכָּרֵת כּל בָּשָׂר עוֹד מִמֵּי הַמַּבּוּל וְלֹא יִהְיֶה עוֹד מַבּוּל לְשַׁחֵת הָאָרֶץ׃}
}
{
	\eP{
		And I shall establish my covenant with you, and all flesh\\
		will not be cut off again by the floodwaters, and there\\
		will not be a flood again to destroy the earth.”\\
	}
}
{
}

\Verse{12}
{
	\hP{וַיֹּאמֶר אֱלֹהִים זֹאת אוֹת הַבְּרִית אֲשֶׁר אֲנִי נֹתֵן בֵּינִי וּבֵינֵיכֶם וּבֵין כּל נֶפֶשׁ חַיָּה אֲשֶׁר אִתְּכֶם לְדֹרֹת עוֹלָם׃}
}
{
	\eP{
		And God said, “This is the sign of the covenant that I’m\\
		giving between me and you and every living being that is\\
		with you for eternal generations:\\
	}
}
{
}

\Verse{13}
{
	\hP{אֶת קַשְׁתִּי נָתַתִּי בֶּעָנָן וְהָיְתָה לְאוֹת בְּרִית בֵּינִי וּבֵין הָאָרֶץ׃}
}
{
	\eP{
		I’ve put my rainbow in the clouds, and it will become a\\
		covenant sign between me and the earth.\\
	}
}
{
%	\footnote{
%		rainbow. It is called this (in Hebrew and in English) because its shape
%		is like a bow that one uses to shoot arrows. The rainbow is the sign of
%		the covenant that promises that God will not destroy the earth again. I
%		heard Rabbi Martin Lawson say that the rainbow symbolizes this because
%		it is “a bow pointed away from the earth.”
%	}
}

\Verse{14}
{
	\hP{וְהָיָה בְּעַנְנִי עָנָן עַל הָאָרֶץ וְנִרְאֲתָה הַקֶּשֶׁת בֶּעָנָן׃}
}
{
	\eP{
		And it will be when I bring a cloud over the earth, and\\
		the rainbow will appear in the cloud,\\
	}
}
{
}

\Verse{15}
{
	\hP{וְזָכַרְתִּי אֶת בְּרִיתִי אֲשֶׁר בֵּינִי וּבֵינֵיכֶם וּבֵין כּל נֶפֶשׁ חַיָּה בְּכל בָּשָׂר וְלֹא יִהְיֶה עוֹד הַמַּיִם לְמַבּוּל לְשַׁחֵת כּל בָּשָׂר׃}
}
{
	\eP{
		and I’ll remember my covenant that is between me and you\\
		and every living being of all flesh, and the waters will\\
		not become a flood to destroy all flesh again.\\
	}
}
{
}

\Verse{16}
{
	\hP{וְהָיְתָה הַקֶּשֶׁת בֶּעָנָן וּרְאִיתִיהָ לִזְכֹּר בְּרִית עוֹלָם בֵּין אֱלֹהִים וּבֵין כּל נֶפֶשׁ חַיָּה בְּכל בָּשָׂר אֲשֶׁר עַל הָאָרֶץ׃}
}
{
	\eP{
		And the rainbow will be in the cloud, and I’ll see it, to\\
		remember an eternal covenant between God and every living\\
		being of all flesh that is on the earth.”\\
	}
}
{
}

\Verse{17}
{
	\hP{וַיֹּאמֶר אֱלֹהִים אֶל נֹחַ זֹאת אוֹת הַבְּרִית אֲשֶׁר הֲקִמֹתִי בֵּינִי וּבֵין כּל בָּשָׂר אֲשֶׁר עַל הָאָרֶץ׃}
}
{
	\eP{
		And God said to Noah, “This is the sign of the covenant\\
		that I’ve established between me and all flesh that is on\\
		the earth.”\\
	}
}
{
}

\Verse{18}
{
	\hJ{וַיִּהְיוּ בְנֵי נֹחַ הַיֹּצְאִים מִן הַתֵּבָה שֵׁם וְחָם וָיָפֶת וְחָם הוּא אֲבִי כְנָעַן׃}
}
{
	\eJ{
		And Noah’s sons who went out from the ark were Shem and\\
		Ham and Yaphet. And Ham: he was the father of Canaan.\\
	}
}
{
}

\Verse{19}
{
	\hJ{שְׁלֹשָׁה אֵלֶּה בְּנֵי נֹחַ וּמֵאֵלֶּה נָפְצָה כל הָאָרֶץ׃}
}
{
	\eJ{
		These three were Noah’s sons, and all the earth expanded\\
		from these.\\
	}
}
{
}

\Verse{20}
{
	\hJ{וַיָּחֶל נֹחַ אִישׁ הָאֲדָמָה וַיִּטַּע כָּרֶם׃}
}
{
	\eJ{
		And Noah began to be a man of the ground, and he planted a\\
		vineyard.\\
	}
}
{
}

\Verse{21}
{
	\hJ{וַיֵּשְׁתְּ מִן הַיַּיִן וַיִּשְׁכָּר וַיִּתְגַּל בְּתוֹךְ אהֳלֹה׃}
}
{
	\eJ{
		And he drank from the wine and was drunk. And he was\\
		exposed inside his tent.\\
	}
}
{
}

\Verse{22}
{
	\hJ{וַיַּרְא חָם אֲבִי כְנַעַן אֵת עֶרְוַת אָבִיו וַיַּגֵּד לִשְׁנֵי אֶחָיו בַּחוּץ׃}
}
{
	\eJ{
		And Ham, the father of Canaan, saw his father’s nakedness\\
		and told his two brothers outside.\\
	}
}
{
}

\Verse{23}
{
	\hJ{וַיִּקַּח שֵׁם וָיֶפֶת אֶת הַשִּׂמְלָה וַיָּשִׂימוּ עַל שְׁכֶם שְׁנֵיהֶם וַיֵּלְכוּ אֲחֹרַנִּית וַיְכַסּוּ אֵת עֶרְוַת אֲבִיהֶם וּפְנֵיהֶם אֲחֹרַנִּית וְעֶרְוַת אֲבִיהֶם לֹא רָאוּ׃}
}
{
	\eJ{
		And Shem and Yaphet took a garment and put it on both\\
		their shoulders and went backwards and covered their\\
		father’s nakedness. And they faced backwards and did not\\
		see their father’s nakedness.\\
	}
}
{
}

\Verse{24}
{
	\hJ{וַיִּיקֶץ נֹחַ מִיֵּינוֹ וַיֵּדַע אֵת אֲשֶׁר עָשָׂה לוֹ בְּנוֹ הַקָּטָן׃}
}
{
	\eJ{
		And Noah woke up from his wine, and he knew what his\\
		youngest son had done to him.\\
	}
}
{
}

\Verse{25}
{
	\hJ{וַיֹּאמֶר אָרוּר כְּנָעַן עֶבֶד עֲבָדִים יִהְיֶה לְאֶחָיו׃}
}
{
	\eJ{
		And he said, “Canaan is cursed. He’ll be a servant of\\
		servants to his brothers.”\\
	}
}
{
}

\Verse{26}
{
	\hJ{וַיֹּאמֶר בָּרוּךְ יְהֹוָה אֱלֹהֵי שֵׁם וִיהִי כְנַעַן עֶבֶד לָמוֹ׃}
}
{
	\eJ{
		And he said, “Blessed is YHWH, God of Shem, and may Canaan\\
		be a servant to them.\\
	}
}
{
}

\Verse{27}
{
	\hJ{יַפְתְּ אֱלֹהִים לְיֶפֶת וְיִשְׁכֹּן בְּאהֳלֵי שֵׁם וִיהִי כְנַעַן עֶבֶד לָמוֹ׃}
}
{
	\eJ{
		May God enlarge Yaphet, and may he dwell in the tents of\\
		Shem, and may Canaan be a servant to them.”\\
	}
}
{
}

\Verse{28}
{
	\hBookOfRecords{וַיְחִי נֹחַ אַחַר הַמַּבּוּל שְׁלֹשׁ מֵאוֹת שָׁנָה וַחֲמִשִּׁים שָׁנָה׃}
}
{
	\eBookOfRecords{
		And Noah lived after the flood three hundred years and\\
		fifty years,\\
	}
}
{
}

\Verse{29}
{
	\hBookOfRecords{וַיִּהְיוּ כּל יְמֵי נֹחַ תְּשַׁע מֵאוֹת שָׁנָה וַחֲמִשִּׁים שָׁנָה וַיָּמֹת׃}
}
{
	\eBookOfRecords{
		and all of Noah’s days were nine hundred years and fifty\\
		years. And he died.\\
	}
}
{
}
